\chapter*{Abstract}
\addcontentsline{toc}{chapter}{Abstract}

Sub-Saharan microgrids frequently operate with sparse, irregularly sampled telemetry, which
undermines classical load-forecasting pipelines that assume dense historical records. This
thesis develops \textbf{Naledi-Forecast}, a hybrid architecture that couples a lightweight
temporal convolutional encoder with a Bayesian correction layer to produce calibrated
short-horizon (15-minute to 6-hour) demand forecasts under data scarcity. The method is
evaluated against three baselines --- seasonal ARIMA, gradient-boosted trees, and a standard
LSTM --- across eighteen simulated microgrid feeders modelled on peri-urban South African
load profiles, plus a twelve-week field deployment on a single university test feeder.

\begin{keybox}{Key Findings}
\begin{itemize}
  \item Naledi-Forecast reduces mean absolute percentage error (MAPE) by \textbf{27.4\%}
        relative to the LSTM baseline under 40\% data-availability conditions.
  \item Calibrated prediction intervals retain \textbf{94.1\%} empirical coverage at the
        nominal 95\% level, versus 81.3\% for uncorrected LSTM intervals.
  \item Inference latency of \textbf{11.8\,ms} per feeder-step makes the model viable for
        embedded controller deployment on constrained hardware.
\end{itemize}
\end{keybox}

The results suggest that Bayesian post-hoc calibration, rather than larger sequence models,
is the more data-efficient lever for improving forecast reliability in scarce-data regimes.
Implications for microgrid dispatch policy and rural electrification planning are discussed
in Chapter~6.

\vspace{0.6cm}
\noindent\textbf{Keywords:} load forecasting, microgrids, Bayesian calibration, temporal
convolutional networks, energy access, data scarcity.

% ============================================================
% ACKNOWLEDGEMENTS
% ============================================================
